% ---------------------------------------------------------------- selected projects
\section{Selected projects}\label{sec:projects}

Case studies first, then shorter entries. Ownership uses five fixed words: \emph{solo} (all
Kartik's work), \emph{lead} (majority author), \emph{co-lead} (a major share), \emph{maintainer}
(took over an existing system) and \emph{contributor} (a minority share). The one team project
says \emph{team of 3} instead. Source links appear where the repository is public.

% ............................................................ GhostShare
\entry[proj:ghostshare]{GhostShare}{Solo · 2026}
\dek{End-to-end-encrypted, ephemeral, real-time sharing: anonymous rooms, no accounts, and a server that only ever sees ciphertext.}

\begin{facts}
\fact{Stack}{TypeScript (strict), Next.js 15, React 19, Fastify 5, Socket.IO 4 with the Redis adapter,
  PostgreSQL 16 and Drizzle, BullMQ, Redis 7, Cloudflare R2, Web Crypto AES-256-GCM, Argon2id
  (hash-wasm), pnpm and Turborepo, Docker Compose behind Coolify, Pino, Sentry}
\fact{Links}{Live at \href{https://share.jkartik.in}{share.jkartik.in} \,·\, write-up:
  \href{https://jkartik.in/blog/zero-knowledge-sharing/}{``The server can't read this''} \,·\,
  \href{https://jkartik.in/projects/ghostshare/}{project page} \,·\, source is private}
\end{facts}

\runin{Problem} Most ``private'' sharing tools still hold plaintext and promise not to read it.
GhostShare is designed so that the server cannot.

\runin{Threat model} The design assumes the server is hostile: the database leaks, object storage
is public, and the operator reads the logs. Each room's secret lives in the URL fragment, which
browsers do not send to the server. The browser derives an AES-256 key from that secret with
Argon2id, a memory-hard function that makes offline guessing of a weak secret expensive, and seals
every note and file with AES-256-GCM, whose authentication tag also exposes tampering by a hostile
server.

\begin{center}
\vspace{2pt}
\begin{tikzpicture}[
  font=\sffamily\footnotesize,
  box/.style={draw=hair, line width=0.5pt, rounded corners=2pt, fill=wash, text=ink,
              align=center, inner xsep=3pt, inner ysep=4pt, minimum width=27mm, minimum height=10mm},
  arrow/.style={-{Stealth[length=4pt,width=3.5pt]}, draw=muted, line width=0.55pt},
  tag/.style={font=\sffamily\scriptsize, text=muted},
  zone/.style={font=\trackedfont\scriptsize, text=muted}]
  \node[box] (sec) at (0,0)   {Room secret in\\URL fragment};
  \node[box] (kdf) at (3.3,0) {Argon2id\\key derivation};
  \node[box] (enc) at (6.6,0) {AES-256-GCM\\encryption};
  \node[box, minimum width=31mm] (api) at (13.3,0.8)  {Fastify + Socket.IO\\ciphertext and room id};
  \node[box, minimum width=31mm] (r2)  at (13.3,-0.8) {Cloudflare R2\\encrypted 5 MiB chunks};
  \draw[arrow] (sec) -- (kdf);
  \draw[arrow] (kdf) -- (enc);
  \draw[arrow] (enc.east) -- (9.3,0) |- (api.west);
  \draw[arrow] (enc.east) -- (9.3,0) |- (r2.west);
  \node[tag, above] at (10.5,0.8) {notes, sync};
  \node[tag, below] at (10.5,-0.8) {signed upload};
  \draw[accent, dashed, line width=0.7pt] (8.65,1.55) -- (8.65,-1.55);
  \node[zone, anchor=south] at (3.3,0.62) {BROWSER: KEY AND PLAINTEXT STAY HERE};
  \node[zone, anchor=south] at (13.3,1.32) {SERVER: ASSUMED HOSTILE};
\end{tikzpicture}\par
\vspace{3pt}
{\footnotesize\itshape\color{muted}Encryption happens before anything reaches the network; the dashed line is the trust boundary.}
\end{center}

\runin{Large files} Files are split into 5 MiB chunks, encrypted in the browser and uploaded straight
to Cloudflare R2 with four parallel uploads, exponential-backoff retry and resume. The app server
hands out short-lived signed upload URLs and never touches the bytes.

\runin{Real time and expiry} Collaboration runs over Socket.IO with a Redis adapter, so it scales
horizontally while relaying only ciphertext. Edits made offline queue in IndexedDB and drain on
reconnect. Every room has a lifetime chosen at creation (1 day, 3 days, 1 week or 1 month), and an
hourly job purges expired rooms and orphaned storage. Helmet and rate limiting guard the API, and
backup and restore scripts cover the data that does persist.

\runin{Limits, stated in the write-up} Metadata such as timing and blob sizes stays visible to the
server, the client code that is served must be trusted, and a room keeps one key for its lifetime;
per-room key rotation is the planned next step.

% ............................................................ NeuroTechh (MAL 4)
\entry[proj:neurotechh]{NeuroTechh: clinic appointment platform}{Team of 3 · Nov 2026}
\dek{The first-prize entry at the MAL 4 hackathon: an architecture and prototype in which patients request appointment slots and clinics confirm or decline them.}

\begin{facts}
\fact{Event}{MAL 4 (Mobile Architecture Lab), Paytm, Mumbai, 1 November 2026. First prize, against working
  professionals and students.}
\fact{Stack}{React Native, Expo, Express, TypeScript, PostgreSQL, Redis, Measure SDK, Docker Compose}
\fact{Links}{Source and design documents:
  \href{https://github.com/KartikJain14/mal-hackathon}{github.com/KartikJain14/mal-hackathon} \,·\, write-up:
  \href{https://jkartik.in/blog/mal4-hackathon-where-code-was-optional/}{``We won a hackathon where the code was optional''}}
\end{facts}

\runin{What was judged} Architecture was 60\% of the score and code was optional. The team of three
submitted a one-page problem brief, a full platform design, an HLD, an LLD, 19 Architecture Decision
Records and three walkthroughs: the core path, failures and a new requirement. The rule the design
protects is that a confirmed slot belongs to exactly one patient and that a cancellation reopens the
slot consistently.

\runin{Key decisions} One modular API rather than the nine microservices first proposed. Postgres row
locks and a partial unique index rather than Redis locks. A transactional outbox rather than Kafka.
Pending requests hold no slot, so the race happens when a clinic confirms rather than when a patient
asks. Idempotency keys are created on the phone.

\runin{Prototype} Built for extra credit: patient, hospital and doctor flows, QR check-in, family
profiles and record uploads. Integration tests cover booking and workflows. The booking rule was
checked with 2, 50 and 200 concurrent confirms on one slot, and exactly one was confirmed each time.

\runin{Limits, stated in the repository} This is a hackathon design and prototype, not a deployed
service. The prototype writes outbox rows in the same transaction as each booking change, but the
worker that would deliver them exists only in the design. The repository's
\href{https://github.com/KartikJain14/mal-hackathon/blob/main/docs/architecture/as-built.md}{as-built
document} records what the code does and where it differs from the design.

% ............................................................ Railway Concession
\entry[proj:concession]{Railway Concession}{Maintainer · Jul 2026 -- present}
\dek{MPSTME's official online railway-concession system: student submissions, collection-slot booking and tools for the college office.}

\begin{facts}
\fact{Ownership}{The original 2025 version was built by the NeuroTechh team. Kartik took it over in
  July 2026 and now maintains it solo, working with the college administration as its rules change.}
\fact{Scale}{Used by \tld4,000 students}
\fact{Stack}{Next.js, PostgreSQL, Drizzle ORM, Nodemailer}
\fact{Links}{Live at \href{https://concession.mpst.me}{concession.mpst.me} \,·\, write-up:
  \href{https://jkartik.in/blog/when-a-form-becomes-infrastructure/}{``When a form becomes infrastructure''}}
\end{facts}

\runin{Data layer} Migrated the system to PostgreSQL with Drizzle ORM.

\runin{Booking rules on the server} A collection date is not accepted just because the browser's date
picker allowed it. One shared module validates every date against the lead time, weekends,
holidays, the configured end date and a rolling window of open dates, computed in Indian Standard
Time whatever the server's own timezone. The client asks the server for the current window only to
guide the student; the submission endpoint validates again before writing. The window slides forward
on its own, so opening a new day needs no deployment.

\runin{Tools for the office} Search by SAP ID with the student's full profile and submission history;
mark a form collected for any date; block SAP IDs; manage holidays; tune the booking window; and
export to Excel. Students get confirmation emails, and the office gets an automated daily
collection-report email. Kartik also patched vulnerable dependencies and improved SEO.

\runin{Outcome} A process that \tld4,000 students depend on, with rules the office can change from a
settings panel rather than waiting on a code change.

% ............................................................ Event & payment backends
\entry[proj:events]{Event and payment backends}{2024 -- 2026}
\dek{Four backends for NMIMS student events, each with one decision that has to be right: whether money changed hands, or whether a badge was already used.}

{\small
\renewcommand{\arraystretch}{1.25}%
\noindent\begin{tabularx}{\linewidth}{@{}>{\raggedright\arraybackslash}p{33mm}>{\raggedright\arraybackslash}p{30mm}>{\raggedright\arraybackslash}p{24mm}>{\raggedright\arraybackslash}X@{}}
\toprule[0.4pt]
\lbl{System} & \lbl{Event} & \lbl{Ownership} & \lbl{Stack and source} \\
\midrule[0.4pt]
Credits and loyalty & Taqneeq 17.0 & Lead author & Flask, Firestore, JWT, Vercel \,·\, \href{https://github.com/KartikJain14/qrshooter-backend}{qrshooter-backend} \\
Delego backend & Mumbai MUN, 2024--25 editions & Lead backend author & FastAPI, PostgreSQL, JWT \,·\, \href{https://github.com/munsoc-mpstme/mundra}{mundra} \\
Paid registration & Programming For Everyone & Co-led & Next.js/Node, PostgreSQL, Cashfree \,·\, \href{https://github.com/MPSTME-ACM/pfe-regi}{pfe-regi} \\
conc-qr & Social Conclave 2026 & Solo & TypeScript, Express, Next.js, PostgreSQL, Sequelize, Docker \,·\, \href{https://github.com/Social-Conclave/conc-qr}{conc-qr} \\
\bottomrule[0.4pt]
\end{tabularx}}

\minientry{Taqneeq 17.0 credits and loyalty backend}{Lead author · 2025}
Taqneeq, the NMIMS MPSTME tech fest, draws thousands of attendees and runs a fest-wide game in which
they earn and spend credits. The backend provides phone-OTP login; a role-based (\code{SALES} /
\code{ADMIN}) credits ledger that records every transaction; referral codes with abuse handling; an
admin portal; and a leaderboard cached with a short TTL, so a crowd refreshing at once reads the
cache instead of the database.

\minientry{Delego backend, Mumbai MUN}{Lead backend author · 2024--25 editions}
The FastAPI and PostgreSQL service behind Delego, the delegate app for Mumbai MUN: passwordless
sign-in with JWT magic links and one-time codes, a QR code per delegate for check-in and meals,
organizer CSV import and export, rate limiting for the morning check-in rush, and scheduled database
backups. The app shipped on the App Store and Play Store for 600+ delegates; those listings have
since been taken down.

\minientry{Paid event registration}{Co-led · 2025}
Paid registration for Programming For Everyone. A registration is confirmed only when Cashfree's
HMAC-verified webhook says the payment cleared, never on a browser redirect, so a closed tab cannot
leave someone charged without a ticket. QR tickets are emailed on confirmation, and gate check-ins
are idempotent, so a re-scan cannot admit someone twice. An admin sales dashboard, Google Sheets
sync, magic-link login and role-based access complete it.

\minientry{conc-qr}{Solo · used 25 Jan 2026}
QR access control for 500+ delegates across Social Conclave's checkpoints, built solo in under two
weeks. Whether a scan is new is decided by a PostgreSQL \code{UNIQUE(participant, access\_type, day)}
constraint: the endpoint attempts the insert and reads the result, so two simultaneous scans cannot
both succeed. Which roles may scan at which checkpoint lives in an access matrix (day \arr
checkpoint \arr roles) that admins edit mid-event without a redeploy. Staff scan with their phone
cameras, with manual entry as a fallback for damaged badges. Write-up:
\href{https://jkartik.in/blog/what-conc-qr-taught-me-about-myself/}{``What conc-qr taught me about myself''}.

% ............................................................ Fool The LLM
\entry[proj:ftl]{Fool The LLM}{Solo · 2026}
\dek{A prompt-injection contest platform: each level is a Gemini model guarding a secret codeword, and players try to talk it into leaking.}

\begin{facts}
\fact{Event}{Started by Kartik as ACM Webmaster; run by ACM and GDG MPSTME. \tld70 participants in
  the most recent edition.}
\fact{Stack}{TypeScript, Next.js, Express, Prisma, PostgreSQL, Gemini}
\fact{Links}{Write-up: \href{https://jkartik.in/blog/where-llm-guardrails-break/}{``Where LLM
  guardrails break''} \,·\, source is private}
\end{facts}

\runin{Design} Levels wrap the model in progressively tighter instructions. The model's own safety
filters are deliberately off, because the guardrail prompt is the game, and that is the surface real
prompt injection attacks. Responses stream token by token; an organizer console manages levels and
participants, shows chat logs live, and can force-unlock or reset a stuck level. Sign-in uses JWT
and magic links, and per-player token rate limits stop one player from spending the event's model
budget.

\runin{What it taught} The write-up draws the defensive lessons: keep secrets out of the model's
context, treat model output as untrusted input, and put the model behind an authorization boundary
rather than asking it to be one.

% ............................................................ bad-db
\entry[proj:baddb]{bad-db}{Solo · 2026}
\dek{A database storage engine written from scratch in Java, to understand what happens between INSERT and the disk.}

\begin{facts}
\fact{Stack}{Java, \code{RandomAccessFile}, a hand-written B-Tree, its own query language (BQL)}
\fact{Links}{Source: \href{https://github.com/KartikJain14/bad-db}{github.com/KartikJain14/bad-db}
  \,·\, write-up: \href{https://jkartik.in/blog/btree-storage-engine/}{``A B-Tree on disk''}}
\end{facts}

Records live in a binary file addressed by byte offset. A hand-written B-Tree, with node splitting
and \code{insertNonFull}, maps primary keys to offsets and is rebuilt from the file on open, so the
file stays the source of truth. Insert, read, update, delete and upsert are supported; deletes write
tombstones, and compaction copies live records into a fresh file to reclaim the space. Older file
formats upgrade automatically. On top sits BQL, a small query language with its own tokenizer,
parser and interactive CLI. The write-up is explicit that the engine is single-threaded and has no
write-ahead log yet.

% ............................................................ ACM infrastructure
\entry[proj:infra]{Self-hosted infrastructure for ACM MPSTME}{Solo · 2025 -- present}
\dek{Infrastructure Kartik set up and runs for the ACM student chapter and other college committees, first as Webmaster and now as Chairperson.}

\begin{itemize}
  \item \textbf{Platform.} Coolify as a self-hosted deployment platform. It hosts
        \href{https://mpstmeacm.com}{mpstmeacm.com}, \href{https://concession.mpst.me}{concession.mpst.me}
        and \href{https://share.jkartik.in}{share.jkartik.in}.
  \item \textbf{Mail.} A Mailcow mail server with SPF, DKIM and DMARC configured by Kartik, shared by
        7--9 college committees with good deliverability. Every ACM member gets an
        \code{@mpstmeacm.com} mailbox on it, at no licence cost.
  \item \textbf{Code.} The chapter's GitHub organisation, GitHub Apps and CI/CD.
  \item \textbf{Community.} An ACM Discord server with custom bots that onboard members and grant
        channel access by position and department.
  \item \textbf{Members-only Minecraft server (2026--27)}, gated by Discord linking through the
        ACM-only bot.
\end{itemize}

% ............................................................ More projects
\subsection{More projects}

\minientry{sports-live: tournament auction and scoring platform}{Solo · freelance · Aug -- Oct 2026}
Used at ELEV8; covered with the freelance work on \pref{proj:sports}.

\minientry[proj:meven]{meven: event and facility requisition portal}{Solo · in development}
A portal for MPSTME that replaces a paper approval chain (faculty mentor \arr institutional
approver \arr admin allotment \arr Joint Registrar escalation) with numbered permission documents
and a full audit trail. One application covers rooms, refuge areas, canopies, guests and
participants, each decided separately. A pnpm monorepo: Next.js, an Express 5 API with a background
worker, PostgreSQL, Redis for queues and rate limiting, Cloudflare R2, SMTP, and Docker with Coolify
on Ubuntu.

\minientry[proj:groupay]{GrouPay: multi-tenant billing for housing societies}{Solo · 2025}
Three tiers of tenancy (super-admin \arr channel admin \arr residents) with the tenancy checks in
middleware, so one society cannot read another's ledger. Invoices with public invoice links, an
audit log, and OTP plus password sign-in. About 15,000 lines in a pnpm monorepo with 8 migrations;
Next.js, Express, Drizzle and PostgreSQL.
Source: \href{https://github.com/KartikJain14/groupay}{github.com/KartikJain14/groupay}.

\minientry[proj:ballast]{ballast: treasury engine for a digital bank}{Solo · hackathon build · Oct 2026}
A ``reserve juggler'' built during a hackathon. Each hour it decides how much cash a bank moving
\tld\$20M an hour should keep liquid and how much to sweep into yield tiers, with a cost-ordered
fallback funding waterfall. It replays 5 policies and runs stress tests, including an SVB-style
bank run; 52 tests check that no decision looks ahead and that the ledger reconciles. Python, with a
web app, a Streamlit console and optional Gemini summaries.
Source: \href{https://github.com/KartikJain14/ballast}{github.com/KartikJain14/ballast}.

\minientry[proj:flap]{flap: an arcade game played over SSH}{Solo · 2026}
A Flappy-style game that runs in the terminal: \code{ssh jkartik.in} and play. Written in Go. The
game state lives on the server, so clients cannot submit scores; the terminal receives a cell-diffed
stream of updates, sprites are half-block pixel art, and a leaderboard ranks players.
Source: \href{https://github.com/KartikJain14/flap}{github.com/KartikJain14/flap}.

\minientry[proj:ai]{captr, usb-ai and hindcribe}{Solo · 2024 -- 2026}
\textbf{captr} (2025) turns a topic into a research document: Gemini writes YouTube search queries,
yt-dlp and transcripts locate the highlights, FFmpeg pulls key frames at those timestamps, and the
result is a Markdown document the user can chat with (FastAPI, MongoDB, JWT and bcrypt, Resend,
Next.js). \textbf{usb-ai} (2026) is a fully offline coding assistant in one Windows executable: a
quantized Gemma-2-2B model embedded in a PyInstaller one-file build and run with llama.cpp, with a
streaming Markdown terminal UI and no network code at all
(\href{https://github.com/KartikJain14/usb-ai}{source}). \textbf{hindcribe} (2024) does offline
Hindi speech recognition and Hindi-to-English translation for the DARPG 2024 public-grievances
hackathon, where citizen audio could not go to a cloud API: FFmpeg, Vosk, Argos Translate and Flask,
with unit tests for each stage (\href{https://github.com/KartikJain14/hindcribe}{source}).

\minientry[proj:bots]{Discord automation bots}{Solo · 2024 -- 2026}
A Taqneeq bot that onboarded 1,100+ hackathon participants with unique invites, real-name renaming
and automatic team channels; an ACM onboarding bot with automatic roles and acceptance emails; a
custom bot for the MUN Society; and the ACM server bots that grant channel access by position and
department.

\minientry[proj:site]{jkartik.in build and deploy pipeline}{Solo · 2026}
An Astro static site shipped as a multi-stage Docker image (Node build \arr nginx-alpine,
\tld97 MB). GitHub Actions with a buildx cache pushes \code{:sha} and \code{:latest} tags to GHCR,
and Watchtower pulls new images to deploy, so CI holds no SSH secrets. nginx with certbot TLS, an
RFC 9116 \code{security.txt}, JSON-LD, and Lighthouse scores of 96 / 100 / 100 / 100.
Source: \href{https://github.com/KartikJain14/jkartik.in}{github.com/KartikJain14/jkartik.in}.

\minientry{Smaller contributions}{}
An adaptive AI department-classifier quiz for Sattva, Taqneeq and Social Conclave, whose questions
adapt to earlier answers (solo or lead). Contributed to a certificate portal: Docker and Coolify
deployment, email sending and fixes. Set up and deployed mpstme.pics: repository scaffold, Docker,
CI/CD with GitHub Actions and Watchtower, and several access-control endpoints.
\href{https://github.com/KartikJain14/template-portfolio}{template-portfolio} (lead), a fill-one-file
portfolio template used by PFE participants.
